\section{System Model}

\subsection{Signal and Impairment Model}

As shown in Fig.~\ref{fig:system-overview}, we consider a coherent FSO downlink from a satellite to an optical ground station. The transmitted symbols use \((8,8)\) 16APSK, with eight equally spaced phase points on each of two amplitude rings, following the modulation definition used by Du \emph{et al.}~\cite{du2025nda}. The constellation is normalized so that its average symbol energy is \(E_s=1\).

Let \(b(k)=\lfloor k/N_{\mathrm{ch}}\rfloor\) identify the atmospheric channel block containing symbol \(k\), where \(N_{\mathrm{ch}}=100\) symbols. The discrete-time complex-baseband observation is
\begin{equation}
  r_k = \sqrt{h_{b(k)}}\,s_k\exp\!\left(j\phi_k\right)+n_k,
  \label{eq:received-signal}
\end{equation}
where \(s_k\) is the transmitted symbol, \(h_{b(k)}\) is the normalized irradiance, and \(n_k\sim\mathcal{CN}(0,N_0)\) is circular complex Gaussian noise. With \(E_s=1\), the average pre-fading data-symbol SNR is \(\bar{\gamma}=E_s/N_0\), its instantaneous value is \(h_{b(k)}\bar{\gamma}\), and \(\gamma_{\mathrm{data}}^{\mathrm{dB}}=10\log_{10}(E_s/N_0)\).

The irradiance follows the Gamma--Gamma model of Al-Habash \emph{et al.}~\cite{alhabash2001} and is generated as the product of two unit-mean Gamma variates. For a plane wave, the small- and large-scale shape parameters are functions of the Rytov variance \(\sigma_R^2\):
\begin{equation}
\begin{aligned}
\alpha&=\left[\exp\!\left(\frac{0.49\sigma_R^2}{(1+1.11(\sigma_R^2)^{6/5})^{7/6}}\right)-1\right]^{-1},\\
\beta&=\left[\exp\!\left(\frac{0.51\sigma_R^2}{(1+0.69(\sigma_R^2)^{6/5})^{5/6}}\right)-1\right]^{-1}.
\end{aligned}
\label{eq:gg-shape-mapping}
\end{equation}
Gu \emph{et al.} use the equivalent plane-wave log-irradiance-variance notation \(\sigma_l^2\) and report weak, moderate, and strong conditions \(\sigma_l^2=0.2,1.6,3.5\) with \((\alpha,\beta)=(11.6,10.1)\), \((4.0,1.9)\), and \((4.2,1.4)\), respectively~\cite{gu2022mfi}. We use those rounded pairs as the three downlink regimes. Substitution into the Al-Habash mapping gives \((11.651,10.122)\), \((4.027,1.911)\), and \((4.226,1.362)\), whose largest componentwise rounding deviation is 2.78\%. The same variance levels trace to the optical-wireless channel-modeling reference cited by Gu \emph{et al.}~\cite{ghassemlooy2019owc}.

Within each evaluated 256-sample processing window, the carrier phase evolves continuously and contains a residual frequency offset \(f_{\mathrm{res}}\), a linear Doppler rate \(\dot f\), and laser phase noise:
\begin{equation}
  \begin{aligned}
  \phi_k
    &= \phi_{\mathrm{init}}+2\pi f_{\mathrm{res}}kT_s
       +\pi\dot f\,(kT_s)^2+\sum_{\ell=0}^{k}\epsilon_\ell,\\
  \epsilon_\ell &\overset{\mathrm{i.i.d.}}{\sim}
       \mathcal{N}\!\left(0,\,2\pi\Delta\nu T_s\right),
  \end{aligned}
  \label{eq:phase-process}
\end{equation}
where \(T_s\) is the symbol interval, \(\Delta\nu=10\)~kHz is the combined laser linewidth, and \(\phi_{\mathrm{init}}\) is the initial phase. The increment variance follows the coherent-FSO linewidth model~\cite{du2025nda}. Each 256-sample window uses an independent realization with continuous phase evolution within the window.

\subsection{Block Structure, Pilot Pattern, and SNR Definition}

Two distinct partitions are used. The atmospheric channel block has \(N_{\mathrm{ch}}=100\) symbols, whereas carrier recovery and adaptive control use nonoverlapping DSP windows of \(N_{\mathrm{DSP}}=256\) samples. A DSP window may therefore span multiple channel blocks; gain estimation follows the 100-symbol partition, while recovery and control use the 256-sample partition.

Within each DSP window, the data-aided layout assigns local indices \(0,4,\ldots,252\) to 64 known pilots and the remaining 192 positions to data. Channel-block boundaries are applied independently of this local layout, so pilots contributing to a channel-gain estimate are selected by their global symbol indices rather than by treating the DSP window as one fading interval. The same transmitted symbols, atmospheric gains, phase sequence, and noise realization are supplied to all compared receiver branches.

For the data-aided branch, one known pilot is placed every \(L_p=4\) transmitted symbols. Thus pilots occupy \(\rho_p=1/4\) of transmitted positions, while the expansion relative to the three data symbols is \(1/3\). Comparisons use a total-transmit-energy coordinate: sending three data symbols plus one pilot requires \(4/3\) of the energy used by three data-only symbols at the same per-symbol energy. The coordinate offset is therefore
\begin{equation}
  \begin{aligned}
  \rho_p&=\frac{1}{4},\\
  \Delta_p&=10\log_{10}\!\left(\frac{4}{3}\right)=1.249~\mathrm{dB},\\
  \gamma_{\mathrm{tot,DA}}^{\mathrm{dB}}
    &=\gamma_{\mathrm{data}}^{\mathrm{dB}}+\Delta_p.
  \end{aligned}
  \label{eq:pilot-energy-coordinate}
\end{equation}
The non-data-aided branch uses every symbol. BER is evaluated over information-bearing positions, with pilot positions excluded from the DA denominator.
